\documentclass[runningheads]{llncs}
\usepackage[utf8]{inputenc}
\usepackage[T1]{fontenc}
\usepackage[polish]{babel}
\usepackage{hyperref}
\usepackage{tikz}
\usetikzlibrary{positioning, arrows.meta, fit}

\tikzset{
  pipeblock/.style={
    rectangle, draw, rounded corners=2pt,
    minimum height=0.95cm, minimum width=1.6cm,
    font=\scriptsize, align=center, inner sep=3pt
  },
  pipearrow/.style={->, >=Stealth, thick}
}

\title{Analiza cross-domain klasyfikacji sentymentu tekstu dla modeli SVM + TF-IDF, LSTM, MiniLM i DistilBERT}
\titlerunning{Analiza cross-domain klasyfikacji sentymentu}
\author{Selega Kamil (288780),\\ Sielska Małgorzata (266560),\\ Szczepaniak Katarzyna (261895),\\ Szeląg Magdalena (266857)}
\authorrunning{K. Selega, M. Sielska, K. Szczepaniak, M. Szeląg}
\institute{
Politechnika Wrocławska\\
 Wydział Informatyki i Telekomunikacji (W4N)\\
 Kierunek: Zaufane Systemy Sztucznej Inteligencji (TAI) \\
 Kurs: Przetwarzanie Języka Naturalnego
}
\date{}

\begin{document}
\maketitle

\section*{Zarys tematu}
Projekt dotyczy analizy odporności modeli klasyfikacji tekstu na zmianę domeny danych (z ang. \textit{out-of-domain text classification}, OOD). Głównym celem ewaluacji jest sprawdzenie, w jakim stopniu modele uczone na jednym typie tekstów potrafią poprawnie klasyfikować dane pochodzące z zupełnie innego źródła. Spadek skuteczności w docelowym środowisku jest powszechnym problemem, dlatego celem pracy jest mierzalna, statystyczna weryfikacja utraty zdolności generalizacyjnych (przy użyciu metryk \textit{accuracy} i \textit{F1}) na zadaniu analizy sentymentu.

\section*{Studium literaturowe}
Problem generalizacji modeli na dane spoza domeny treningowej od dawna stanowi jedno z większych wyzwań przetwarzania języka naturalnego. Ramy teoretyczne tego zjawiska opisują m.in. Yang i in. \cite{yang2023outofdistributiongeneralizationtextclassification} oraz Lang i in. \cite{lang2023surveyoutofdistributiondetectionnlp}. Zwracają oni uwagę, że pomimo świetnych wyników współczesnych sieci neuronowych w obrębie znanych danych docelowych (in-domain), skuteczność klasyfikacji drastycznie spada przy zetknięciu z innym kontekstem wypowiedzi. Początki prób przeciwdziałania temu zjawisku sięgają klasycznej adaptacji domenowej nakierowanej jeszcze na stare klasyfikatory statystyczne, co analizowali chociażby Daumé III i Marcu \cite{DBLP:journals/corr/abs-1109-6341}. Wraz z narastającym skomplikowaniem sieci, uwaga analityków przesunęła się m.in. na rozpoznawanie nowych intencji z ujęciem metod nienadzorowanych (bez etykiet nowej domeny) \cite{9741317} czy ugruntowanie pewności i bezpieczeństwa najnowszych zautomatyzowanych chatbotów konwersacyjnych \cite{9052492}.

Trudności z generalizacją na dane spoza domeny wymuszają nowe metody analiz w zróżnicowanych gałęziach NLP: obniżają jakość tłumaczenia maszynowego \cite{haddow-koehn-2012-analysing}, ograniczają ekstrakcję encji (ang. \textit{named entity recognition}) podczas analiz tekstów medycznych \cite{9648338}, a nawet generują szum decyzyjny przy segmentowaniu izolowanych tekstów języków azjatyckich przeplatanych wtrąceniami europejskimi \cite{Limkonchotiwat2021HandlingCA}. Zjawisko wykracza daleko poza ramy NLP i dotyczy też kategoryzacji obrazów wizyjnych, chociażby klasyfikacji środowiska na podstawie rzutów satelitarnych \cite{Kolodner2022EvaluatingDL}.

Obok głębokich ujęć rekurencyjnych (LSTM) standardami w zadaniach kategoryzacji zdań i wyszukiwania tekstu stały się od niedawna reprezentacje oparte natywnie na mechanizmie uwagi (Transformer, m.in. popularny BERT). Szerokiej weryfikacji takich architektur podjęli się badawczo i polecają takie metody m.in. Sharma \cite{Sharma2025ACA} oraz Oday i in. \cite{Oday2025TextCF}. Istotną luką pozostaje jednak ograniczanie rozmiaru tych modeli. Stąd naukowcy zwrócili się ku mniejszym, ale zoptymalizowanym przez transfer learning wersjom tych samych modeli, na czele z DistilBERT-em. Barbon i Akabane \cite{s22218184} potwierdzili badawczo wysoką, często zaskakującą skuteczność lekkiego DistilBERT-a w niwelowaniu różnic między domenami w analizach wielojęzycznych w porównaniu do starszych architektur klasyfikacyjnych; podobne wnioski przedstawili Muffo i Bertino, testując zjawisko na dużych zasobach tekstowych dla języka włoskiego \cite{muffo2023bertinoitaliandistilbertmodel}. Co najważniejsze z perspektywy badanego w tym projekcie problemu, Yuan \cite{Yuan2026EmpiricalSO} potwierdził na zbiorach IMDb skuteczność lekkiej wersji DistilBERT, wykazując przewagę klasyfikacji opartej na Transformerze nad klasycznymi sieciami głębokimi (CNN i LSTM). Jednocześnie warto wspomnieć, że badając tak złożone architektury należy posługiwać się decyzjami popartymi jasnymi wyjaśnieniami (ang. \textit{explanations}), co dla modeli OOD szerzej wykazali Chrysostomou i Aletras \cite{chrysostomou2022empiricalstudyexplanationsoutofdomain}.

\section*{Plan eksperymentu}
Opierając się na powyższym przeglądzie, zdecydowano się przetestować algorytmy na zadaniu klasyfikacji sentymentu stosując ewaluację międzydomenową (ang. \textit{cross-domain}). Zamiast ustalać wyłącznie jedną domenę bazową, każda z zademonstrowanych architektur będzie trenowana niezależnie dla każdego poszczególnego zbioru danych. Po wytrenowaniu modelu na danym zbiorze (źródłowym), zostanie on przetestowany na nim samym (ewaluacja \textit{in-domain}, dalej IND) oraz na każdym z pozostałych (ewaluacja \textit{out-of-domain}, dalej OOD). Badanie rotacyjnie obejmie poniższe zbiory:
\begin{itemize}
    \item \textbf{IMDb Dataset}: 50\,000 recenzji filmowych, zbalansowane (25\,000 / 25\,000); etykiety z pola \textit{sentiment}. Często używane jako bazowy punkt odniesienia jakości modeli głębokich w analizie sentymentu \cite{Yuan2026EmpiricalSO,imdb_dataset}.
    \item \textbf{Rotten Tomatoes Dataset}: spośród oryginalnego zbioru wylosowano 50\,000 zrównoważonych recenzji krytyków filmowych (25\,000 / 25\,000); etykiety z pola \textit{review\_type}: \textit{fresh}$\rightarrow$1, \textit{rotten}$\rightarrow$0 \cite{rotten_tomatoes}.
    \item \textbf{Amazon Reviews Dataset}: 21\,055 recenzji produktów, niezbalansowane (6\,705 pozytywnych / 14\,350 negatywnych); etykiety z pola \textit{Rating} (ocena $\geq 3$ oznacza recenzję pozytywną) \cite{amazon_reviews}.
    \item \textbf{Yelp Reviews}: 10\,000 recenzji lokalnych usług, niezbalansowane (8\,324 pozytywne / 1\,676 negatywnych); etykiety z pola \textit{stars} (ocena $\geq 3$ oznacza recenzję pozytywną) \cite{yelp_dataset}.
\end{itemize}
Wszystkie teksty przetworzono wspólnym potokiem czyszczącym: konwersja do małych liter, usunięcie znaczników HTML, adresów URL, znaków niealfabetycznych oraz normalizacja białych znaków. Niezbalansowanie zbiorów Amazon i Yelp zostało świadomie zachowane, co stanowi materiał do trzeciego pytania badawczego.
Dzięki zbudowaniu pełnej macierzy testowej (każdy model trenowany kolejno na jednej z czterech domen i testowany na wszystkich czterech) badanie pozwoli zmierzyć, w jakim stopniu zmiana domeny wpływa na parametry dokładności modeli, których tło teoretyczne opisują Yang i Lang \cite{yang2023outofdistributiongeneralizationtextclassification,lang2023surveyoutofdistributiondetectionnlp}. W ramach protokołów testowych przyjmuje się rygorystyczny podział każdego ze środowisk danych. Aby zapewnić wiarygodność doboru hiperparametrów i uniknąć przetrenowania na zbiorze źródłowym, na etapie treningu zastosowana zostanie $k$-krotna walidacja krzyżowa (ang. \textit{k-fold cross-validation}). Finalna ewaluacja każdego modelu przeprowadzona zostanie natomiast na wydzielonych zbiorach testowych, zarówno IND, jak i OOD, co umożliwi miarodajne i statystycznie poprawne porównania cross-domain.

Do weryfikacji celów postawiono na zróżnicowany stack algorytmiczny, aby porównać skuteczność pomiędzy starymi standardami klasyfikacyjnymi i innowacyjnym użyciem uwag i transformatorów \cite{DBLP:journals/corr/abs-1109-6341,s22218184}. Posłużą temu cztery stopnie trudności badanej architektury:
\begin{enumerate}
    \item \textbf{SVM + wektoryzacja TF-IDF}, reprezentująca optymalny wczesny model z nurtu klasyfikatorów statystycznych, tworząca próg bazowego startu eksperymentu.
    \item \textbf{LSTM}, stanowiący stabilniejszy, bazowy standard przetwarzania, nierzadko słabszy od modeli Transformerowych na danych spoza domeny, co Yuan \cite{Yuan2026EmpiricalSO} wskazuje jako punkt odniesienia do pomiaru siły tych różnic.
    \item \textbf{MiniLM (sentence-transformers)}, czyli lekki transformator zdaniowy używany w trybie zamrożonym jako ekstraktor cech, stanowiący kompromis pomiędzy klasycznym LSTM a pełnym dostrajaniem modeli językowych \cite{wang2020minilm,reimers-2019-sentence-bert}. Pierwotnie planowany w jego miejscu pełny BERT okazał się zbyt kosztowny obliczeniowo w protokole 5-fold $\times$ 4 domeny (szacowany czas wykonania całego protokołu na lokalnym GPU przekraczał 18 godzin), dlatego rolę silnego transformatorowego \textit{baseline'u} przejął MiniLM.
    \item \textbf{DistilBERT}, będący wariantem pozwalającym zbadać, czy skompresowany model po transfer learningu zachowuje odporność na zmianę domeny porównywalną z pełnym BERT-em (teza zaczerpnięta z pracy Barbona \cite{s22218184,sanh2019distilbert}).
\end{enumerate}

\section*{Pytania badawcze}
W badaniu sformułowano trzy pytania badawcze, z których każde dotyczy innej właściwości metod lub danych i jest weryfikowane odrębnym podzbiorem zaplanowanych eksperymentów:
\begin{enumerate}
    \item Jak duży jest spadek skuteczności klasyfikacji sentymentu po zmianie domeny (OOD) względem oceny IND? Pytanie dotyczy właściwości \emph{danych} (dystans rozkładów źródłowy--docelowy). Eksperyment: wyliczenie różnic miar \textit{accuracy} oraz \textit{F1 macro} pomiędzy oceną IND a OOD dla każdej komórki pełnej macierzy 4$\times$4.
    \item Czy modele oparte na uwadze (MiniLM, DistilBERT) generalizują lepiej między domenami niż klasyczne podejścia (SVM + TF-IDF, LSTM)? Pytanie dotyczy właściwości \emph{metod} (indukcyjne biasy architektur). Eksperyment: porównanie średnich spadków OOD-vs-IND oraz absolutnych wartości \textit{F1 macro} pomiędzy czterema badanymi architekturami.
    \item W jakim stopniu charakterystyka domeny źródłowej (rozmiar zbioru, balans klas) wpływa na zdolność generalizacyjną modelu do pozostałych domen? Pytanie łączy właściwości \emph{danych} i \emph{metod}. Eksperyment: analiza wierszowa macierzy cross-domain, czyli porównanie średniego \textit{F1 macro} OOD osiąganego po treningu na każdej z czterech domen, ze szczególnym uwzględnieniem zbiorów niezbalansowanych (Amazon, Yelp) względem zbalansowanych (IMDb, Rotten Tomatoes).
\end{enumerate}

\section*{Opis zaimplementowanych metod}
Cztery zaimplementowane modele dobrano tak, by pokryć stopniowy gradient skomplikowania reprezentacji języka, od dyskretnej reprezentacji \textit{bag-of-words} aż po pełne dostrajanie modelu Transformerowego. Każdy model został zaimplementowany w niezależnym notatniku Jupyter z izolowanymi zależnościami.

\subsection*{SVM + TF-IDF}
Klasyczny statystyczny model bazowy. Reprezentację tekstu stanowi \textit{TfidfVectorizer} ze \textit{scikit-learn} z bigramami (\textit{ngram\_range}=(1,2)), z włączonym \textit{sublinear\_tf} i ograniczeniem do 50\,000 cech. Klasyfikatorem jest \textit{LinearSVC} ($C$=1.0, \textit{class\_weight="balanced"}, do 5\,000 iteracji, \textit{dual="auto"}). Wektoryzator jest dopasowywany niezależnie w każdym foldzie do tekstów treningowych, a na zbiorach testowych OOD stosowana jest jedynie transformacja, co odzwierciedla realistyczny scenariusz wdrożenia.

\begin{figure}[h]
\centering
\begin{tikzpicture}[node distance=0.45cm]
  \node[pipeblock] (n1) {Tekst\\surowy};
  \node[pipeblock, right=of n1] (n2) {Czyszczenie\\(lower, regex)};
  \node[pipeblock, right=of n2] (n3) {TF-IDF\\(1--2)-gramy\\50\,000 cech};
  \node[pipeblock, right=of n3] (n4) {LinearSVC\\$C{=}1.0$\\balanced};
  \node[pipeblock, right=of n4] (n5) {Etykieta\\$\{0,1\}$};
  \draw[pipearrow] (n1) -- (n2);
  \draw[pipearrow] (n2) -- (n3);
  \draw[pipearrow] (n3) -- (n4);
  \draw[pipearrow] (n4) -- (n5);
\end{tikzpicture}
\caption{Potok klasyfikacji SVM + TF-IDF.}
\end{figure}

\subsection*{LSTM}
Rekurencyjna sieć neuronowa zaimplementowana w \textit{TensorFlow/Keras}. Architektura: warstwa \textit{Embedding} (rozmiar słownika 20\,000, wymiar wektora 128) $\rightarrow$ \textit{LSTM} (128 jednostek ukrytych) $\rightarrow$ \textit{Dropout} (0,5) $\rightarrow$ \textit{Dense}(1, sigmoid). Tokenizator \textit{Keras Tokenizer} jest wspólny dla wszystkich domen, długość sekwencji obcinana jest do 95. percentyla długości w danej domenie. Optymalizacja: \textit{Adam}, funkcja straty \textit{binary crossentropy}, rozmiar batcha 64, do 20 epok treningu z mechanizmem \textit{EarlyStopping}.

\begin{figure}[h]
\centering
\resizebox{\textwidth}{!}{%
\begin{tikzpicture}[node distance=0.4cm]
  \node[pipeblock] (n1) {Tekst\\surowy};
  \node[pipeblock, right=of n1] (n2) {Tokenizer\\Keras\\($|V|{=}20\,000$)};
  \node[pipeblock, right=of n2] (n3) {Padding\\do p95\\długości};
  \node[pipeblock, right=of n3] (n4) {Embedding\\$20\,000\times 128$};
  \node[pipeblock, right=of n4] (n5) {LSTM\\128 jedn.};
  \node[pipeblock, right=of n5] (n6) {Dropout\\$p{=}0{,}5$};
  \node[pipeblock, right=of n6] (n7) {Dense\\$\sigma$};
  \node[pipeblock, right=of n7] (n8) {$P(y{=}1)$};
  \draw[pipearrow] (n1) -- (n2);
  \draw[pipearrow] (n2) -- (n3);
  \draw[pipearrow] (n3) -- (n4);
  \draw[pipearrow] (n4) -- (n5);
  \draw[pipearrow] (n5) -- (n6);
  \draw[pipearrow] (n6) -- (n7);
  \draw[pipearrow] (n7) -- (n8);
\end{tikzpicture}%
}
\caption{Potok klasyfikacji LSTM (kolejność warstw od wejścia tekstowego do prawdopodobieństwa klasy pozytywnej).}
\end{figure}

\subsection*{MiniLM (zamrożony enkoder + regresja logistyczna)}
Lekki Transformer zdaniowy w trybie ekstraktora cech. Wykorzystano model \textit{sentence-transformers/all-MiniLM-L6-v2} \cite{wang2020minilm,reimers-2019-sentence-bert} z \textit{max\_seq\_length}=256, zwracający 384-wymiarowe znormalizowane embeddingi zdań. Na ich podstawie trenowana jest \textit{LogisticRegression} (\textit{liblinear}, \textit{class\_weight="balanced"}, do 5\,000 iteracji). Embeddingi są obliczane raz dla każdej domeny i reużywane we wszystkich foldach, co znacząco redukuje czas eksperymentu względem fine-tuningu Transformera \textit{end-to-end}.

\begin{figure}[h]
\centering
\begin{tikzpicture}[node distance=0.45cm]
  \node[pipeblock] (n1) {Tekst\\surowy};
  \node[pipeblock, right=of n1] (n2) {Tokenizer\\WordPiece\\($L_{\max}{=}256$)};
  \node[pipeblock, right=of n2, minimum width=2.2cm] (n3) {MiniLM L6\\(zamrożony)\\$\rightarrow$ embedding $384$-d};
  \node[pipeblock, right=of n3] (n4) {Logistic\\Regression\\(\,liblinear, balanced\,)};
  \node[pipeblock, right=of n4] (n5) {Etykieta\\$\{0,1\}$};
  \draw[pipearrow] (n1) -- (n2);
  \draw[pipearrow] (n2) -- (n3);
  \draw[pipearrow] (n3) -- (n4);
  \draw[pipearrow] (n4) -- (n5);
  \node[below=0.15cm of n3, font=\scriptsize\itshape] {wagi zamrożone};
\end{tikzpicture}
\caption{Potok klasyfikacji MiniLM: zamrożony enkoder zdaniowy dostarcza embeddingi, na których trenowana jest regresja logistyczna.}
\end{figure}

\subsection*{DistilBERT (fine-tuning)}
Pełne dostrojenie zdestylowanej wersji modelu BERT (\textit{distilbert-base-uncased}, 6 warstw, $\sim$66M parametrów) \cite{sanh2019distilbert} z biblioteki \textit{Hugging Face Transformers}. Wejście tokenizowane do długości 128. Optymalizacja: \textit{AdamW} z $\eta=2\times10^{-5}$, rozmiar batcha 32, do 3 epok z \textit{early stopping} o cierpliwości 2.
\begin{figure}[h]
\centering
\begin{tikzpicture}[node distance=0.45cm]
  \node[pipeblock] (n1) {Tekst\\surowy};
  \node[pipeblock, right=of n1] (n2) {Tokenizer\\WordPiece\\($L_{\max}{=}128$)};
  \node[pipeblock, right=of n2, minimum width=2.2cm] (n3) {DistilBERT\\(fine-tuning)\\6 warstw, $\sim$66M};
  \node[pipeblock, right=of n3] (n4) {Głowica\\liniowa\\(2 wyjścia)};
  \node[pipeblock, right=of n4] (n5) {Etykieta\\$\{0,1\}$};
  \draw[pipearrow] (n1) -- (n2);
  \draw[pipearrow] (n2) -- (n3);
  \draw[pipearrow] (n3) -- (n4);
  \draw[pipearrow] (n4) -- (n5);
  \node[below=0.15cm of n3, font=\scriptsize\itshape] {wszystkie wagi aktualizowane (AdamW, $\eta{=}2\!\times\!10^{-5}$)};
\end{tikzpicture}
\caption{Potok klasyfikacji DistilBERT z pełnym dostrajaniem (\textit{end-to-end}) - w przeciwieństwie do MiniLM aktualizowane są wszystkie wagi enkodera.}
\end{figure}

\section*{Protokół ewaluacji cross-domain}
Każdy z czterech modeli był trenowany niezależnie dla każdej z czterech domen, dając pełną macierz 4$\times$4 par (domena treningowa, domena testowa). Stosowana jest 5-krotna walidacja krzyżowa (\textit{StratifiedKFold} dla SVM, MiniLM i DistilBERT; \textit{KFold} dla LSTM, ze względu na uwarunkowania implementacyjne potoku treningowego), z ustalonym ziarnem losowości \textit{random\_state}=42 dla powtarzalności.

W każdym z 5 foldów, po wytrenowaniu modelu na części treningowej domeny źródłowej, przeprowadzane są dwie ewaluacje:
\begin{itemize}
    \item \textbf{In-domain (IND)}: na wydzielonej części walidacyjnej tej samej domeny;
    \item \textbf{Out-of-domain (OOD)}: na całych pozostałych trzech zbiorach (przy zachowaniu wektoryzatora / tokenizatora / wag modelu dopasowanych do domeny źródłowej).
\end{itemize}
Daje to 80 punktów pomiarowych na każdy z czterech modeli (4 domeny treningowe $\times$ [1 IND + 3 OOD] domeny testowe $\times$ 5 foldów). Mierzone metryki to: \textit{accuracy}, \textit{precision}, \textit{recall} oraz \textit{F1} (uśrednienie \textit{macro}, dające równą wagę klasie pozytywnej i negatywnej, co jest istotne wobec niezbalansowania zbiorów Amazon i Yelp). Zachowanie wyników na poziomie pojedynczego folda umożliwi podczas dalszej analizy zbadanie istotności statystycznej zaobserwowanych różnic.

\section*{Środowisko badawcze i artefakty}
Eksperymenty zaimplementowano w Pythonie 3.13 (zarządzanie zależnościami przez \textit{uv}), w czterech niezależnych notatnikach Jupyter (po jednym na model), w celu izolacji łańcuchów zależności pomiędzy klasycznym \textit{stackiem} \textit{scikit-learn}, \textit{TensorFlow/Keras} oraz \textit{PyTorch + Hugging Face}. Wykorzystane biblioteki: \textit{scikit-learn} (SVM, regresja logistyczna, walidacja krzyżowa i metryki), \textit{TensorFlow/Keras} (LSTM), \textit{Hugging Face Transformers} z \textit{PyTorch} (DistilBERT), \textit{sentence-transformers} (MiniLM). Wszystkie przetworzone zbiory danych, wytrenowane modele LSTM, dzienniki treningu (\textit{training\_logs.csv}) oraz pełne wyniki cross-domain (\textit{ood\_results.csv}) są wersjonowane w repozytorium projektu, co zapewnia pełną powtarzalność eksperymentów na potrzeby dalszych analiz.

\bibliographystyle{plain}
\bibliography{bibliography}

\end{document}
